\section{Introduction}
\label{sec:intro}

SFT distillation of an RL-trained agent can appear to succeed by format-level tool-call metrics---and still fail catastrophically in deployment. We distill MEM1~\cite{zhou2025mem1}, a Qwen2.5-7B retrieval-augmented QA agent trained with PPO via veRL~\cite{sheng2024verl}, into a Qwen2.5-3B student using 1{,}118 teacher trajectories, and we replicate the entire pipeline with a Llama-3.2-1B student (cross-family, ${\sim}1/3$ the Qwen-3B student size) to test whether the failure is backbone-specific or systematic (\S\ref{sec:cross_backbone}). On both backbones, format-level tool-call indicators rate the student at $1.00$ (tool-call syntax, balanced XML tags, non-empty queries; Qwen query Jaccard $=0.95$), while content-level metrics reveal the collapse: BLEU $=1.14/1.16$ on sacrebleu's 0--100 scale, ROUGE-L $=0.21/0.21$, and em\_partial at $N\!=\!8$ retaining only $0.021/0.014$ of teacher capability (full per-metric breakdown in Appendix~\ref{app:silent_failure_full}). Action-sequence accuracy is $\approx 0.606$ for both students, but per-sample forensics (Appendix~\ref{app:asa_forensics}) reveals this coincidence is structurally degenerate---both students emit the identical action prefix---rather than an independent diagnostic signal. These content warnings are easy to dismiss as cosmetic noise because the student produces teacher-like \emph{structure}---a blind spot mirroring the ``syntax-only'' tool-call evaluations critiqued by \citet{jhandi2025slm} (\S\ref{sec:related}).

It did not succeed on either backbone. The Qwen-3B student achieves $\text{em\_partial} = 0.021$ at $N\!=\!8$ simultaneous objectives---a $95.4\%$ drop from the teacher's $0.455$---with 83\% of outputs truncated by hallucinated search traces. The Llama-1B student reaches $\text{em\_partial} = 0.014$ at $N\!=\!8$ with 86\% truncation---a near-identical collapse despite the change in both family and parameter count. The failure is \emph{silent} at the format/syntax layer on both backbones: the format-level checks all pass while content-level metrics (BLEU, ROUGE-L, em\_partial) expose the collapse\footnote{Action-sequence accuracy $\approx 0.606$ for both students, but this is a structurally degenerate metric: both backbones emit the identical 3-token action prefix on all 125 test samples (per-sample forensics in Appendix~\ref{app:asa_forensics}).}. The collapse surfaces only when search calls must return real results, i.e., at pipeline deployment. This is the central danger: an invisible failure---reproducible across backbones---that passes format-level tool-call quality gates while content-level warning signs are overridden by teacher-like structure.

Why do richer training signals not help? At the matched-sample scale available (176 pairs), DPO does not close the gap ($+1.0$pp, $p = 0.163$). Scaling from 3B to 7B SFT closes only ${\sim}12$--$14\%$ of the teacher--student gap (\S\ref{sec:capacity}), ruling out capacity as the primary bottleneck.

This paper is deliberately scoped as a \emph{case study} on MEM1-family ReAct agents---a single teacher, a single multi-hop QA domain, and two student backbones---rather than a universal claim about agentic distillation. An earlier attempt to cross-validate the protocol on Kang et al.'s CodeAct-family checkpoints~\cite{kang2025agent} (exp\_015) was abandoned at the Phase~0 dry-run stage because positional oracle-replay injection presupposes a ReAct-style \texttt{<search>}/\texttt{<information>} schema and does not transfer to free-form Python-action protocols; the honest boundary is that our empirical validation is restricted to the ReAct family (\S\ref{sec:related} Kang paragraph, Appendix~\ref{app:redaction} Deviation log).

We introduce a diagnostic protocol that distinguishes \emph{strategy acquisition} from \emph{execution capability}. The key causal test: providing oracle retrieval (referred to throughout as ``oracle'' for brevity; technically a \textit{teacher-retrieval replay}---we replay the teacher's actual retrieval output without any answer-key oracle) at inference recovers em\_partial from $0.021$ to $0.498$ at $N\!=\!8$ ($+47.7$pp, $p < 0.0001$), with truncation dropping from 83\% to 0\%. This gain is a \emph{leakage-inclusive upper bound}: a passage-redaction ablation (\S\ref{sec:causal}, Appendix~\ref{app:redaction}) yields a \emph{partial-mechanism} decomposition in which retrieval quality alone contributes a \emph{structural retrieval residual} of ${\sim}11$pp (${\sim}23\%$; $+11.0$pp after redacting gold-answer substrings from oracle passages, gain measured over the \texttt{honest\_search} FlashRAG baseline), while passage-level answer-leakage-assisted decoding contributes the remaining $+36.7$\,pp $[+33.5, +40.0]$ ($\approx 77\%$; Table~\ref{tab:redaction_dose_response}, Appendix~\ref{app:redaction}). Both components are causally meaningful---retrieval is \emph{necessary but not sufficient}---and together establish that the student has learned a viable search strategy but cannot exploit \emph{non-substring-hint} passages at the level its teacher can. Even a simple BM25 retrieval system recovers $+6$--$12$pp (within the teacher's training document pool). The protocol proceeds by elimination---collapse characterization and weight geometry analysis (\S\ref{sec:framework}), causal intervention (\S\ref{sec:distillation})---and localizes the failure before deployment (Figure~\ref{fig:framework}).

\begin{figure}[t]
\centering
\includegraphics[width=\columnwidth]{framework_overview.pdf}
\caption{Overview of the \textsc{MemDistill} diagnostic framework. Phase~1 (top) progressively characterizes the distillation failure through collapse analysis, SVD weight geometry, and activation probing. Phase~2 (middle) tests interventions: DPO preference learning, truncation confound correction, and oracle retrieval. The capacity bound (at most ${\sim}12$--$14\%$) is an independent upper bound; the leakage-inclusive oracle recovery is further decomposed (passage-redaction, exp\_016) into retrieval-quality ($\sim 11$\,pp, $\sim 23\%$) and passage-level answer-leakage ($\sim 37$\,pp, $\sim 77\%$).}
\label{fig:framework}
\end{figure}

\paragraph{Contributions.}
\begin{itemize}[leftmargin=*,nosep]
    \item \textbf{A silent failure mode in agentic SFT distillation}: format-level indicators (tool-call syntax 1.00, balanced XML 1.00, query Jaccard 0.95) give high scores while task performance retains <6\% of teacher capability at N=8.
    \item A \textbf{diagnostic protocol} (collapse characterization $\to$ weight geometry $\to$ causal intervention) that localizes the failure to execution-environment mismatch before pipeline deployment. We propose and demonstrate on MEM1 a three-step elimination scheme (capacity $\to$ evaluation artefacts $\to$ causal intervention) for diagnosing silent-failure-shaped gaps.
    \item \textbf{Mechanism decomposition}: oracle (teacher-retrieval replay) $+47.7$pp ($p < 0.0001$) is a \emph{leakage-inclusive} upper bound; a passage-redaction ablation (\S\ref{sec:causal}) decomposes it into retrieval-quality $\sim 11$pp ($23\%$) and passage-level answer-leakage $\sim 37$pp ($77\%$). Capacity upper-bounds the gap at ${\sim}12$--$14\%$, acting as an independent upper bound (not a partition); DPO is an informative negative (power $> 0.99$ at 5pp).
    \item A \textbf{deployable mitigation}: BM25 retrieval recovers $+6$--$12$pp with no additional training, immediately deployable as a sidecar component within the teacher's training document pool.
\end{itemize}
